\documentclass[11pt]{article}

\usepackage[margin=1in]{geometry}
\usepackage{amsmath}
\usepackage{amssymb}
\usepackage{mathtools}
\usepackage{booktabs}
\usepackage{enumitem}
\usepackage{hyperref}

\newcommand{\dd}{\mathrm{d}}
\newcommand{\grad}{\nabla}
\newcommand{\divergence}{\nabla\!\cdot}
\newcommand{\curl}{\nabla\times}
\newcommand{\bnd}{\operatorname{bound}}
\newcommand{\sign}{\operatorname{sign}}
\newcommand{\diag}{\operatorname{diag}}
\newcommand{\vol}{\mathcal{V}}
\newcommand{\faces}{\mathcal{F}}
\newcommand{\nb}{\mathcal{N}}

\title{Mathematical and Numerical Formulation of the Sharp-Interface Free-Surface Dam-Break Model}
\author{MOOSE Navier-Stokes Module}
\date{\today}

\begin{document}

\maketitle

\begin{abstract}
This document describes the mathematics and numerical implementation of the
sharp-interface free-surface dam-break model used by
\texttt{dam\_break\_openfoam\_geometry.i}.  The model is a two-phase,
immiscible, incompressible, one-velocity volume-of-fluid formulation coupled to
the MOOSE Navier-Stokes module through mixture-property functors, a
reduced-pressure finite-volume pressure equation, a conservative
alpha-consistent mass flux, and a Rhie-Chow pressure-velocity coupling.  The
notation and several modeling ideas are intentionally close to common
VOF/PIMPLE descriptions such as OpenFOAM's \texttt{interFoam} documentation,
but the equations and numerical details below describe the MOOSE
implementation used by this submodule.
\end{abstract}

\tableofcontents

\section{Nomenclature}

The nomenclature and principal symbols used in this document are the following:

\begin{itemize}[leftmargin=2.5em]
  \item \(\alpha\) (--) liquid volume fraction
  \item \(\alpha_g=1-\alpha\) (--) gas volume fraction
  \item \(\rho_l,\rho_g\) (kg\,m\(^{-3}\)) liquid and gas densities
  \item \(\rho(\alpha)\) (kg\,m\(^{-3}\)) mixture density
  \item \(\mu_l,\mu_g\) (Pa\,s) liquid and gas dynamic viscosities
  \item \(\mu(\alpha)\) (Pa\,s) mixture dynamic viscosity
  \item \(\mathbf{u}\) (m\,s\(^{-1}\)) mixture velocity
  \item \(p\) (Pa) physical pressure
  \item \(p_{rgh}\) (Pa) reduced pressure, \(p-\rho\mathbf{g}\cdot(\mathbf{x}-\mathbf{x}_{ref})\)
  \item \(\mathbf{g}\) (m\,s\(^{-2}\)) gravitational acceleration
  \item \(gh\) (m\(^2\)\,s\(^{-2}\)) gravitational head, \(\mathbf{g}\cdot(\mathbf{x}-\mathbf{x}_{ref})\)
  \item \(c_\alpha\) (--) interface-compression coefficient
  \item \(\mathbf{n}_{\alpha}\) (--) regularized interface normal
  \item \(\delta_n\) (--) interface-normal regularization
  \item \(P,N\) cell indices; \(P\) is owner/element, \(N\) is neighbor
  \item \(f\) face index
  \item \(V_P\) (m\(^3\)) cell volume
  \item \(A_f\) (m\(^2\)) face area/measure
  \item \(\mathbf{n}_f\) (--) unit face normal in owner-to-neighbor orientation
  \item \(\varphi_f\) (m\,s\(^{-1}\)) volumetric normal flux density
  \item \(F_f=A_f\varphi_f\) (m\(^3\)\,s\(^{-1}\)) integrated volumetric flux
  \item \(F_{\alpha,f}\) (m\(^3\)\,s\(^{-1}\)) integrated volume-fraction flux
  \item \(F_{\rho,f}\) (kg\,s\(^{-1}\)) integrated mass flux
  \item \(F_{H/A,f}\) (m\(^3\)\,s\(^{-1}\)) integrated momentum-predictor flux
  \item \(F_{p,f}\) (m\(^3\)\,s\(^{-1}\)) integrated pressure-equation flux
  \item \(D_f\) face-normal pressure coefficient
  \item \(rAU_P\) inverse momentum diagonal
  \item \(N_s\) number of alpha subcycles
  \item \(\Delta t\), \(\Delta t_s\) global and alpha-subcycle timesteps
  \item \(Co_\alpha\) alpha Courant number
  \item \(\lambda_f\) MULES-style face limiter
  \item \(I_{\mathrm{backflow}}\) backflow indicator at an open boundary
\end{itemize}

\section{Dam-Break Configuration}

The input represents the OpenFOAM-parity dam-break geometry:

\begin{align}
  a &= 0.05715\ \mathrm{m}, \\
  L_x &= 10a, \\
  L_y &= 1.25a,
\end{align}

with a \(400\times 50\) Cartesian finite-volume mesh.  The initial liquid column
is

\begin{equation}
  0 \le x < a,
  \qquad
  0 \le y < a.
\end{equation}

The material constants are

\begin{align}
  \rho_l &= 998.19,
  &
  \rho_g &= 1.185, \\
  \mu_l &= 10^{-3},
  &
  \mu_g &= 1.48\times 10^{-5},
\end{align}

and

\begin{equation}
  \mathbf{g}=(0,-9.81,0).
\end{equation}

The active numerical parameters in the free-surface path are

\begin{align}
  \Delta t &= 10^{-4}, &
  N_{\alpha,\mathrm{sub}} &= 2, \\
  c_\alpha &= 0.01, &
  N_{\alpha,\mathrm{corr}} &= 1, \\
  N_{\lambda,\mathrm{iter}} &= 6.
\end{align}

No surface-tension coefficient, turbulence model, energy equation, or phase
change model is configured in this input.  The model is therefore a laminar,
isothermal, immiscible, incompressible two-phase free-surface calculation.

\section{Continuous Free-Surface Model}

\subsection{Volume Fractions and Mixture Properties}

The free surface is represented by the liquid volume fraction

\begin{equation}
  0 \le \alpha(\mathbf{x},t) \le 1.
\end{equation}

The gas fraction is

\begin{equation}
  \alpha_g = 1-\alpha.
\end{equation}

The mixture laws used by the Navier-Stokes kernels are affine in \(\alpha\):

\begin{align}
  \rho(\alpha)
  &=
  \alpha\rho_l + (1-\alpha)\rho_g
  =
  \rho_g + (\rho_l-\rho_g)\alpha,
  \label{eq:rho-alpha}
  \\
  \mu(\alpha)
  &=
  \alpha\mu_l + (1-\alpha)\mu_g.
  \label{eq:mu-alpha}
\end{align}

These mixture properties are exposed to the Navier-Stokes module as functor
material properties.  Momentum, pressure, and boundary kernels consume the
resulting \(\rho(\alpha)\) and \(\mu(\alpha)\) fields rather than owning a
separate phase model.

\subsection{Reduced Pressure}

Define

\begin{equation}
  gh(\mathbf{x}) = \mathbf{g}\cdot(\mathbf{x}-\mathbf{x}_{ref}).
\end{equation}

The reduced pressure is

\begin{equation}
  p_{rgh}=p-\rho gh.
  \label{eq:prgh}
\end{equation}

Taking a gradient gives

\begin{equation}
  \grad p
  =
  \grad p_{rgh}
  +
  gh\,\grad\rho
  +
  \rho\grad gh.
\end{equation}

Because \(\grad gh=\mathbf{g}\),

\begin{equation}
  -\grad p+\rho\mathbf{g}
  =
  -\grad p_{rgh}-gh\,\grad\rho.
  \label{eq:reduced-pressure-force}
\end{equation}

Equation \eqref{eq:reduced-pressure-force} is the mathematical basis for the
sharp-interface reduced-pressure force.  The pressure equation solves for the
reduced pressure variable, and the density-gradient hydrostatic correction is
handled explicitly through the sharp-interface face-flux machinery.

\subsection{Governing Equations}

The incompressibility constraint is

\begin{equation}
  \divergence\mathbf{u}=0.
  \label{eq:div-free}
\end{equation}

The liquid volume fraction satisfies a conservative VOF equation with an
interface-compression contribution,

\begin{equation}
  \frac{\partial\alpha}{\partial t}
  +
  \divergence(\alpha\mathbf{u})
  +
  \divergence\!\left[
    c_\alpha\alpha(1-\alpha)\mathbf{u}_c
  \right]
  =
  0.
  \label{eq:continuous-alpha}
\end{equation}

In the MOOSE implementation, \(\mathbf{u}_c\) is not an independent field.
Instead, the VOF corrector constructs a compressed face state from the transport
flux direction, the interface normal, and \(c_\alpha\).  This is the same
numerical role played by artificial compression in VOF/PIMPLE solvers such as
OpenFOAM's \texttt{interFoam}: sharpen the smeared interface without introducing
a physical mass flux.

The reduced-pressure momentum equation is

\begin{equation}
  \frac{\partial(\rho\mathbf{u})}{\partial t}
  +
  \divergence(\rho\mathbf{u}\otimes\mathbf{u})
  -
  \divergence
  \left[
    \mu
    \left(\grad\mathbf{u}+\grad\mathbf{u}^{T}\right)
  \right]
  =
  -\grad p_{rgh}
  -
  gh\,\grad\rho.
  \label{eq:momentum-prgh}
\end{equation}

The viscous stress above is written in the simplified incompressible laminar
form used by this dam-break case.  If a future calculation enabled turbulence
or surface tension, those terms would enter the momentum equation through the
standard Navier-Stokes module kernels and materials; they are not active in the
present input.

\section{Initial and Boundary Data}

The initial alpha field is the liquid-column indicator:

\begin{equation}
  \alpha(\mathbf{x},0)
  =
  \begin{cases}
    1, & x<a\ \text{and}\ y<a, \\
    0, & \text{otherwise}.
  \end{cases}
\end{equation}

The initial velocity is

\begin{equation}
  \mathbf{u}(\mathbf{x},0)=\mathbf{0}.
\end{equation}

The initialized reduced pressure is

\begin{equation}
  p_{rgh}(\mathbf{x},0)
  =
  \begin{cases}
    -(\rho_l-\rho_g)g(L_y-a), & x<a\ \text{and}\ y<a, \\
    0, & \text{otherwise}.
  \end{cases}
  \label{eq:init-prgh}
\end{equation}

This is not an absolute hydrostatic pressure field; it is the supplied
reduced-pressure state.  The startup projection described later repairs the
discrete face flux while restoring this pressure vector.

At solid walls,

\begin{equation}
  \mathbf{u}=\mathbf{0}.
\end{equation}

The wall pressure treatment imposes a pressure flux/normal-gradient constraint
consistent with the pressure predictor.  In simplified notation,

\begin{equation}
  \frac{\partial p_{rgh}}{\partial n}
  =
  -\frac{F_{\mathrm{predictor},b}}{D_b},
  \label{eq:wall-pressure-gradient}
\end{equation}

unless the Rhie-Chow object has already populated a constrained wall gradient,
in which case the cached gradient is used directly.

At the top outlet, the pressure condition is a reduced-pressure total-pressure
condition,

\begin{equation}
  p_{rgh,b}
  =
  p_{ref}
  -
  \frac{1}{2}\rho_b
  |\mathbf{u}_{n,b}|^2 I_{\mathrm{backflow}}
  -
  \rho_b gh_b.
  \label{eq:top-prgh}
\end{equation}

The input sets \(p_{ref}=0\).  The alpha boundary condition is inlet-outlet:

\begin{equation}
  \alpha_b =
  \begin{cases}
    0, & F_b<0\quad\text{(backflow into the domain)},\\
    \alpha_P + O(d_{Pb}\grad\alpha_P), & F_b\ge 0\quad\text{(outflow)}.
  \end{cases}
\end{equation}

\section{Interface Geometry}

The flow physics creates geometry functors from the alpha field.  The clipped
alpha field used for interface indicators is

\begin{equation}
  \alpha_c=\min(\alpha_{\max},\max(\alpha_{\min},\alpha)).
\end{equation}

The near-interface indicator is

\begin{equation}
  \chi_I =
  \begin{cases}
    1, & \alpha_{I,lo}\le \alpha_c\le \alpha_{I,hi},\\
    0, & \text{otherwise}.
  \end{cases}
\end{equation}

The regularized interface normal is

\begin{equation}
  \mathbf{n}_{\alpha}
  =
  \frac{\grad\alpha}{|\grad\alpha|+\delta_n}.
  \label{eq:interface-normal}
\end{equation}

For the dam-break case, this normal is consumed by the VOF compression
correction through the face-oriented functor
\texttt{flow\_interface\_unit\_normal\_face}.  The default regularization is
\(\delta_n=10^{-8}\) unless changed by input.

\section{Cell-Centered Finite-Volume Discretization}

\subsection{Control-Volume Balance}

Let \(P\) be a finite-volume cell with volume \(V_P\).  A generic conservation
law

\begin{equation}
  \frac{\partial\psi}{\partial t}
  +
  \divergence\mathbf{j}_{\psi}
  =
  s_\psi
\end{equation}

integrates to

\begin{equation}
  V_P
  \frac{\psi_P^{n+1}-\psi_P^n}{\Delta t}
  +
  \sum_{f\in\partial P}
  J_{\psi,f}
  =
  V_P s_{\psi,P}.
  \label{eq:fv-balance}
\end{equation}

For advective transport,

\begin{equation}
  J_{\psi,f}=F_f\psi_f,
  \qquad
  F_f=A_f\varphi_f.
\end{equation}

The face normal \(\mathbf{n}_f\) is oriented from the owner cell \(P\) to the
neighbor cell \(N\).  Hence \(F_f>0\) is owner-to-neighbor transport.

\subsection{Scalar Transport Matrix}

The donor/upwind part of a scalar advection equation leads to a linear system

\begin{equation}
  a_P\psi_P+\sum_{N\in\nb(P)}a_N\psi_N=b_P.
\end{equation}

For a pure transient upwind advection equation, the owner-cell coefficients can
be interpreted as

\begin{align}
  a_P
  &=
  \frac{V_P}{\Delta t}
  +
  \sum_{f\in\partial P}\max(F_f,0),
  \\
  a_N
  &=
  \min(F_f,0).
\end{align}

Boundary conditions contribute either matrix terms, source terms, or both.  The
VOF corrector does not assemble this matrix; it acts after the upwind matrix
solve by applying a conservative limited correction to the alpha flux.

\subsection{Pressure Diffusion Operator}

On an internal face, the pressure normal derivative can be decomposed into an
orthogonal two-point part plus an explicit non-orthogonal correction:

\begin{equation}
  \left(\frac{\partial p}{\partial n}\right)_f
  =
  \frac{p_N-p_P}{\mathbf{d}_{PN}\cdot\mathbf{n}_f}
  +
  \mathbf{k}_f\cdot(\grad p)_f.
\end{equation}

The dam-break input disables orthogonality correction in the flow physics, so
the dominant contribution is the two-point normal difference.  The pressure
diffusion flux has the abstract form

\begin{equation}
  F_{p,f}
  =
  M_{P,f}p_P+M_{N,f}p_N-R_f,
  \label{eq:pressure-discrete-flux}
\end{equation}

where \(M_{P,f}\) and \(M_{N,f}\) are matrix contributions from the pressure
diffusion kernel, and \(R_f\) is its explicit right-hand-side contribution.

\section{VOF Transport and Subcycling}

\subsection{Alpha Courant Number}

The VOF transport flux is frozen during each alpha solve.  The alpha Courant
number used to decide whether more subcycles are needed is

\begin{equation}
  Co_{\alpha,P}
  =
  \frac{1}{2}
  \frac{\Delta t}{V_P}
  \sum_{f\in\partial P}
  |F_f|.
  \label{eq:alpha-courant}
\end{equation}

The global alpha Courant number is

\begin{equation}
  Co_\alpha=\max_P Co_{\alpha,P}.
\end{equation}

The number of alpha subcycles is

\begin{equation}
  N_s
  =
  \max
  \left(
    N_{\alpha,\mathrm{sub}},
    \left\lceil
      \frac{Co_\alpha}{Co_{\alpha,\max}}
    \right\rceil
  \right),
  \qquad
  \Delta t_s=\frac{\Delta t}{N_s}.
  \label{eq:subcycles}
\end{equation}

For subcycle \(s\),

\begin{align}
  t_{old,s}&=t^n+s\Delta t_s,\\
  t_s&=t^n+(s+1)\Delta t_s.
\end{align}

The old alpha state is advanced between subcycles but restored to the true
timestep-old state after all subcycles.  This preserves the global time
history expected by the transient Navier-Stokes machinery.

\subsection{Donor Solve}

For each subcycle, the base matrix solve computes an intermediate alpha
\(\alpha_P^*\):

\begin{equation}
  \frac{V_P}{\Delta t_s}
  \left(\alpha_P^*-\alpha_P^{old}\right)
  +
  \sum_{f\in\partial P}
  F_f\alpha_{upwind,f}
  =
  0.
  \label{eq:alpha-donor}
\end{equation}

The upwind state is

\begin{equation}
  \alpha_{upwind,f}
  =
  \begin{cases}
    \alpha_P, & F_f\ge 0,\\
    \alpha_N, & F_f<0,
  \end{cases}
  \label{eq:upwind-alpha}
\end{equation}

for internal faces.  Boundary faces use the active finite-volume boundary
condition.

The donor alpha flux is

\begin{equation}
  F_{\alpha,f}^{donor}=F_f\alpha_{upwind,f}.
  \label{eq:donor-flux}
\end{equation}

The donor solve is monotone but diffusive.  The MOOSE sharp-interface VOF path
therefore splits the final flux into

\begin{equation}
  F_{\alpha,f}^{final}
  =
  F_{\alpha,f}^{donor}
  +
  F_{\alpha,f}^{limited,corr}.
  \label{eq:flux-split}
\end{equation}

The correction is a conservative face-flux correction, not a post-solve clip of
cell values.  This distinction is essential: clipping cell alpha would generally
destroy conservation, while a limited face flux preserves pairwise cancellation
on internal faces.

\section{High-Order Compression Flux}

\subsection{High-Order Face State}

For internal faces, the corrector evaluates a bounded Van Leer face value:

\begin{equation}
  \alpha_f^{HO}
  =
  \bnd\left(\alpha_f^{VanLeer}\right),
  \qquad
  \bnd(a)=\min(1,\max(0,a)).
\end{equation}

The linearly interpolated alpha used by the compression term is

\begin{equation}
  \alpha_f^{lin}
  =
  \bnd
  \left[
    g_C\alpha_P+(1-g_C)\alpha_N
  \right],
  \label{eq:linear-alpha-face}
\end{equation}

where \(g_C\) is the geometric interpolation weight stored on the face.

\subsection{Compression State}

Using the regularized interface normal \eqref{eq:interface-normal}, the
compressed alpha face state is

\begin{equation}
  \alpha_f^{comp}
  =
  \bnd
  \left[
    \alpha_f^{HO}
    +
    c_\alpha
    \sign(F_f)
    \alpha_f^{lin}(1-\alpha_f^{lin})
    \left(\mathbf{n}_{\alpha,f}\cdot\mathbf{n}_f\right)
  \right].
  \label{eq:compressed-alpha}
\end{equation}

The target alpha flux is

\begin{equation}
  F_{\alpha,f}^{target}
  =
  F_f\alpha_f^{comp}.
  \label{eq:target-alpha-flux}
\end{equation}

The factor \(\alpha_f^{lin}(1-\alpha_f^{lin})\) localizes compression near the
interface.  It vanishes in pure liquid and pure gas cells.  The sign of \(F_f\)
orients the compression contribution with the transport direction, while
\(\mathbf{n}_{\alpha,f}\cdot\mathbf{n}_f\) projects the interface normal onto
the face normal.

\subsection{Boundary Face Classification}

The VOF corrector distinguishes internal, Dirichlet inflow, Dirichlet outflow,
open outflow, and closed boundary faces.  A compact representation of the donor
state is

\begin{equation}
  \alpha_{donor,f}
  =
  \begin{cases}
    \alpha_P,
      & F_f\ge 0,\\
    \alpha_N,
      & F_f<0,\quad f\ \text{internal},\\
    \alpha_b,
      & F_f<0,\quad f\ \text{Dirichlet inflow},\\
    \alpha_b^{backflow},
      & F_f<0,\quad f\ \text{open boundary backflow},\\
    \alpha_P,
      & f\ \text{closed}.
  \end{cases}
\end{equation}

Closed faces do not contribute a correction flux.  On the top outlet of the
dam-break case, backflow imposes gas, \(\alpha_b^{backflow}=0\), and outflow
uses the extrapolated interior alpha state.

\section{MULES-Style Bounded Flux Limiter}

\subsection{Raw Correction}

Let \(F_{\alpha,f}^{work}\) be the current working alpha flux.  At the start of
the first correction sweep,

\begin{equation}
  F_{\alpha,f}^{work}=F_{\alpha,f}^{donor}.
\end{equation}

The raw correction is

\begin{equation}
  C_f
  =
  F_{\alpha,f}^{corr}
  =
  F_{\alpha,f}^{target}
  -
  F_{\alpha,f}^{work}.
  \label{eq:raw-correction}
\end{equation}

If applied without limiting, the owner-cell update would be

\begin{equation}
  \Delta\alpha_{P,f}
  =
  -\frac{\Delta t_s}{V_P}C_f.
  \label{eq:owner-correction-update}
\end{equation}

Thus \(C_f>0\) removes alpha from the owner cell and \(C_f<0\) adds alpha to
the owner cell.  The neighbor cell receives the opposite signed update.

\subsection{Local Admissible Bounds}

The implementation initializes a hard admissible interval \([0,1]\), then
tightens it with bounded neighbor or boundary values.  Let \(B_P\) be the set
of such adjacent values collected for cell \(P\).  Define

\begin{align}
  \alpha_P^{lo}
  &=
  \max
  \left(
    0,\max_{\beta\in B_P}\bnd(\beta)
  \right),
  \\
  \alpha_P^{hi}
  &=
  \min
  \left(
    1,\min_{\beta\in B_P}\bnd(\beta)
  \right).
\end{align}

The correction capacity that can increase alpha is

\begin{equation}
  \Psi_P^{inc}
  =
  \frac{V_P}{\Delta t_s}
  \max(0,\alpha_P^{lo}-\alpha_P^*),
  \label{eq:psi-inc}
\end{equation}

and the correction capacity that can decrease alpha is

\begin{equation}
  \Psi_P^{dec}
  =
  \frac{V_P}{\Delta t_s}
  \max(0,\alpha_P^*-\alpha_P^{hi}).
  \label{eq:psi-dec}
\end{equation}

These names describe the effect on the cell alpha.  The source code stores the
same quantities as two maps used to limit negative-inflow and positive-outflow
corrections.

\subsection{Cell Flux Sums}

For cell \(P\), write \(C_{P,f}\) as the correction flux signed so that positive
values remove alpha from \(P\).  The raw correction sums are

\begin{align}
  S_P^{out,+}
  &=
  \sum_{f\in\partial P}
  \max(0,C_{P,f}),
  \\
  S_P^{in,-}
  &=
  \sum_{f\in\partial P}
  \max(0,-C_{P,f}).
\end{align}

During limiter iteration \(m\), the already accepted limited sums are

\begin{align}
  S_{P,m}^{out,+}
  &=
  \sum_{f\in\partial P}
  \max(0,\lambda_f^{(m)}C_{P,f}),
  \\
  S_{P,m}^{in,-}
  &=
  \sum_{f\in\partial P}
  \max(0,-\lambda_f^{(m)}C_{P,f}).
\end{align}

The cell-side limiter estimates are

\begin{align}
  \lambda_P^{inc}
  &=
  \bnd
  \left(
    \frac{S_{P,m}^{out,+}+\Psi_P^{inc}}
         {S_P^{in,-}+\epsilon}
  \right),
  \\
  \lambda_P^{dec}
  &=
  \bnd
  \left(
    \frac{S_{P,m}^{in,-}+\Psi_P^{dec}}
         {S_P^{out,+}+\epsilon}
  \right),
  \label{eq:cell-limiters}
\end{align}

where \(\epsilon=\sqrt{\mathrm{minReal}}\) in the implementation.

\subsection{Face Limiter}

For an internal face with owner \(P\), neighbor \(N\), and owner-oriented
correction \(C_f\), the accepted face limiter is tightened as

\begin{equation}
  \lambda_f
  \leftarrow
  \begin{cases}
    \min(\lambda_f,\lambda_P^{dec},\lambda_N^{inc}),
      & C_f>0,\\
    \min(\lambda_f,\lambda_P^{inc},\lambda_N^{dec}),
      & C_f<0.
  \end{cases}
  \label{eq:face-limiter}
\end{equation}

On a boundary face, only the owner-side limiter participates, unless the face is
closed.  On a processor-partition face, both ranks compute a limiter and the
minimum is exchanged so the two sides apply the same \(\lambda_f\).

The limited correction on alpha-correction sweep \(k\) is

\begin{equation}
  F_{\alpha,f}^{limited,corr}
  =
  w_k\lambda_f C_f,
  \qquad
  w_k=
  \begin{cases}
    1, & k=0,\\
    0.5, & k>0.
  \end{cases}
  \label{eq:limited-correction}
\end{equation}

The dam-break input uses \(N_{\alpha,\mathrm{corr}}=1\), so \(w_0=1\).

The owner and neighbor updates are

\begin{align}
  \Delta\alpha_P
  &=
  -\frac{\Delta t_s}{V_P}
  F_{\alpha,f}^{limited,corr},
  \\
  \Delta\alpha_N
  &=
  +\frac{\Delta t_s}{V_N}
  F_{\alpha,f}^{limited,corr}.
\end{align}

Therefore an internal face conserves integrated alpha exactly:

\begin{equation}
  V_P\Delta\alpha_P+V_N\Delta\alpha_N=0.
\end{equation}

After the update, the working flux is advanced:

\begin{equation}
  F_{\alpha,f}^{work}
  \leftarrow
  F_{\alpha,f}^{work}
  +
  F_{\alpha,f}^{limited,corr}.
\end{equation}

\section{Subcycle-Averaged Alpha and Mass Fluxes}

The corrector publishes timestep-consistent subcycle-averaged face fluxes.  Let

\begin{equation}
  \theta_s=\frac{\Delta t_s}{\Delta t}.
\end{equation}

The published alpha flux is

\begin{equation}
  \bar{F}_{\alpha,f}
  =
  \sum_{s=0}^{N_s-1}
  \theta_s F_{\alpha,f,s}^{work}.
  \label{eq:published-alpha-flux}
\end{equation}

The alpha-consistent mass flux for subcycle \(s\) is

\begin{equation}
  F_{\rho,f,s}
  =
  \rho_g F_{f,s}
  +
  (\rho_l-\rho_g)
  F_{\alpha,f,s}^{work}.
  \label{eq:rho-flux-subcycle}
\end{equation}

The published mass flux is

\begin{equation}
  \bar{F}_{\rho,f}
  =
  \sum_{s=0}^{N_s-1}
  \theta_s F_{\rho,f,s}.
  \label{eq:published-rho-flux}
\end{equation}

This flux is exposed to the Navier-Stokes momentum equation as the conservative
sharp-interface mass flux.  It is generally not equal to
\(\rho_f\bar{F}_f\) with an independently interpolated \(\rho_f\).

\section{Alpha-Consistent Density Conservation}

From \eqref{eq:rho-alpha},

\begin{equation}
  V_P(\rho_P^{n+1}-\rho_P^n)
  =
  (\rho_l-\rho_g)
  V_P(\alpha_P^{n+1}-\alpha_P^n).
\end{equation}

Using the conservative alpha update gives

\begin{equation}
  V_P(\rho_P^{n+1}-\rho_P^n)
  =
  -\Delta t
  \sum_{f\in\partial P}
  (\rho_l-\rho_g)\bar{F}_{\alpha,f}.
\end{equation}

The mass flux \eqref{eq:published-rho-flux} satisfies

\begin{equation}
  \sum_{f\in\partial P}\bar{F}_{\rho,f}
  =
  \rho_g
  \sum_{f\in\partial P}\bar{F}_f
  +
  (\rho_l-\rho_g)
  \sum_{f\in\partial P}\bar{F}_{\alpha,f}.
\end{equation}

For a pressure-corrected incompressible transport flux,

\begin{equation}
  \sum_{f\in\partial P}\bar{F}_f\approx 0,
\end{equation}

so

\begin{equation}
  V_P(\rho_P^{n+1}-\rho_P^n)
  +
  \Delta t
  \sum_{f\in\partial P}
  \bar{F}_{\rho,f}
  \approx 0.
  \label{eq:density-consistency}
\end{equation}

Equation \eqref{eq:density-consistency} is the central reason the
free-surface submodule publishes \(\bar{F}_{\rho,f}\).  The momentum equation
should see the same phase transport that updated \(\alpha\), especially for the
large density ratio in the dam-break problem.

\section{Momentum Predictor}

For each velocity component \(i\), the segregated linearized momentum equation
has the abstract form

\begin{equation}
  a_{P,i}u_{P,i}
  +
  \sum_{N\in\nb(P)}
  a_{N,i}u_{N,i}
  =
  b_{P,i}
  -
  V_P(\grad p_{rgh})_{P,i}
  +
  V_P f_{rgh,P,i}.
  \label{eq:momentum-discrete}
\end{equation}

Here

\begin{equation}
  \mathbf{f}_{rgh}=-gh\,\grad\rho
\end{equation}

is the reduced-pressure density-gradient force from
\eqref{eq:reduced-pressure-force}.  Define

\begin{equation}
  rAU_{P,i}=\frac{1}{a_{P,i}}.
\end{equation}

The off-diagonal and explicit source contribution is grouped as \(H_{P,i}\),
giving

\begin{equation}
  u_{P,i}
  =
  HbyA_{P,i}
  -
  rAU_{P,i}(\grad p_{rgh})_{P,i}
  +
  rAU_{P,i}f_{rgh,P,i},
  \qquad
  HbyA_{P,i}=rAU_{P,i}H_{P,i}.
\end{equation}

The pressure equation and Rhie-Chow correction use the same \(rAU\)-based face
coefficient that comes from this momentum operator.  This consistency is more
important than the precise algebraic splitting of \(H\) and \(A\): the pressure
diffusion coefficient must correspond to the momentum predictor it corrects.

\section{Pressure Equation}

\subsection{Predictor and Hydrostatic Fluxes}

The pressure equation is obtained by enforcing continuity on the pressure
corrected face flux.  The predictor contribution is

\begin{equation}
  F_{H/A,f}
  =
  A_f
  \left[
    (\mathbf{HbyA})_f\cdot\mathbf{n}_f
    +
    \varphi_{ddtCorr,f}
  \right],
  \label{eq:fhbya}
\end{equation}

where \(\varphi_{ddtCorr,f}\) is the transient projection correction used when
startup/predictor source terms are enabled.

The explicit reduced-pressure hydrostatic contribution is

\begin{equation}
  F_{g,f}
  =
  A_f
  \left[
    -gh_f
    (\grad_n\rho)_f
    D_f
  \right].
  \label{eq:phig}
\end{equation}

The face-normal pressure coefficient is

\begin{equation}
  D_f=\mathbf{n}_f\cdot rAU_f\mathbf{n}_f.
  \label{eq:pressure-coefficient}
\end{equation}

The pressure-divergence flux functor uses the sign convention

\begin{equation}
  F_{\phi H/A,f}
  =
  -(F_{H/A,f}+F_{g,f}).
  \label{eq:stored-phi-hbya}
\end{equation}

\subsection{Pressure Correction Flux}

After the pressure solve, the pressure diffusion flux is cached in the form
\eqref{eq:pressure-discrete-flux}.  The corrected integrated volumetric flux is

\begin{equation}
  F_f^{corr}
  =
  -F_{\phi H/A,f}
  +
  F_{p,f}.
  \label{eq:corrected-integrated-flux}
\end{equation}

The pressure equation enforces

\begin{equation}
  \sum_{f\in\partial P} F_f^{corr}=0.
  \label{eq:pressure-continuity}
\end{equation}

The face flux density published for transport is

\begin{equation}
  \varphi_f^{corr}=\frac{F_f^{corr}}{A_f}.
\end{equation}

This value is the authoritative pressure-corrected volumetric transport flux.
The next alpha solve freezes it into the VOF transport flux so alpha subcycles
see a fixed velocity field.

\section{Rhie-Chow Pressure-Velocity Coupling}

On a collocated cell-centered mesh, simple interpolation of cell velocities to
faces can decouple pressure and velocity.  A plain interpolated face flux

\begin{equation}
  \varphi_f^{lin}
  =
  \left[
    g_C\mathbf{u}_P+(1-g_C)\mathbf{u}_N
  \right]\cdot\mathbf{n}_f
\end{equation}

does not necessarily contain the nearest-neighbor pressure difference
\((p_N-p_P)\).  The resulting pressure equation can admit an alternating
checkerboard mode.

Rhie-Chow interpolation restores a compact pressure-difference contribution:

\begin{equation}
  \varphi_f^{RC}
  =
  \varphi_f^{lin}
  -
  D_f
  \left[
    \frac{p_N-p_P}{\mathbf{d}_{PN}\cdot\mathbf{n}_f}
    -
    \overline{\grad p}_f\cdot\mathbf{n}_f
  \right]
  +
  \varphi_{\mathrm{src},f}.
  \label{eq:rhie-chow}
\end{equation}

Here \(\varphi_{\mathrm{src},f}\) collects explicit source corrections, including
the sharp-interface hydrostatic contribution when it is enabled.  The bracketed
term vanishes for a pressure field whose cell-to-cell difference and
interpolated pressure gradient are consistent, but it penalizes pressure modes
that are invisible to a simple collocated interpolation.

The sharp-interface implementation keeps the face flux, not the reconstructed
cell velocity, as the source of truth.  The cell-centered pressure-corrected
velocity is reconstructed from face-normal pressure-correction scalars:

\begin{equation}
  \delta\mathbf{u}_P
  \approx
  rAU_P\,
  \mathcal{R}_P
  \left[
    \frac{F_{p,f}+F_{g,f}}{A_fD_f}
  \right],
  \label{eq:velocity-writeback}
\end{equation}

where \(\mathcal{R}_P\) denotes the finite-volume reconstruction from face
normal scalars to a cell vector.  This reconstructed velocity is used for the
next momentum state and boundary evaluations, but it is not projected back to
overwrite \(F_f^{corr}\).

\section{Startup Projection}

At the first timestep, the initial pressure and velocity fields may not produce
a face flux that is discrete-continuity-consistent with the pressure operator
and boundary constraints.  The startup projection solves a pressure-only
continuity correction while preserving the supplied reduced-pressure field.

Let \(p_0\) be the initialized reduced-pressure vector and let \(F_f^0\) be the
initial predictor flux.  The startup projection computes an auxiliary pressure
increment \(\pi\) such that

\begin{equation}
  \sum_{f\in\partial P}
  \left[
    F_f^0
    -
    A_fD_f
    \left(\frac{\partial\pi}{\partial n}\right)_f
  \right]
  =
  0.
  \label{eq:startup-projection}
\end{equation}

The startup-corrected flux is

\begin{equation}
  F_f^{startup}
  =
  F_f^0
  -
  A_fD_f
  \left(\frac{\partial\pi}{\partial n}\right)_f.
\end{equation}

After the projection, the pressure vector is restored:

\begin{equation}
  p_{rgh}^{after\ startup}=p_0.
\end{equation}

The startup guard suppresses the explicit hydrostatic and startup predictor
source terms during this cleanup:

\begin{equation}
  F_{g,f}^{startup}=0,
  \qquad
  F_{ddtCorr,f}^{startup}=0.
\end{equation}

Thus the startup projection repairs continuity of the face flux without
replacing the physical reduced-pressure initialization \eqref{eq:init-prgh}.

\section{Coupling Contracts with the Navier-Stokes Module}

The sharp-interface free-surface submodule does not replace the Navier-Stokes
finite-volume discretization.  Instead, it supplies the mathematical fields
that the standard Navier-Stokes kernels need:

\begin{enumerate}[leftmargin=2.5em]
  \item The scalar field \(\alpha\), from which \(\rho(\alpha)\) and
        \(\mu(\alpha)\) are generated.
  \item The interface normal \(\mathbf{n}_{\alpha}\), used only by the VOF
        compression correction in this case.
  \item The frozen alpha-transport flux \(F_f\), obtained from the previous
        pressure-corrected volumetric face flux.
  \item The bounded alpha flux \(\bar{F}_{\alpha,f}\), produced by the MULES-style
        corrector.
  \item The alpha-consistent mass flux \(\bar{F}_{\rho,f}\), consumed by
        momentum advection.
  \item The reduced-pressure face coefficients \(D_f\) and predictor fluxes
        used by the pressure equation and pressure boundary conditions.
\end{enumerate}

The mathematical coupling order in one timestep is:

\begin{enumerate}[leftmargin=2.5em]
  \item freeze the latest pressure-corrected volumetric flux for alpha
        transport;
  \item solve the donor alpha equation \eqref{eq:alpha-donor};
  \item apply the bounded correction \eqref{eq:limited-correction};
  \item publish \(\bar{F}_{\alpha,f}\) and \(\bar{F}_{\rho,f}\);
  \item update \(\rho(\alpha)\), \(\mu(\alpha)\), and interface geometry;
  \item solve the momentum predictor \eqref{eq:momentum-discrete};
  \item solve the pressure equation \eqref{eq:pressure-continuity};
  \item publish \(F_f^{corr}\) for the next transport state.
\end{enumerate}

This order is the main correctness contract.  The momentum and pressure
equations see material properties and mass fluxes consistent with the alpha
field that was just transported.

\section{Conservation and Boundedness}

For any set of cells \(\Omega_h\), the alpha update satisfies

\begin{equation}
  \sum_{P\in\Omega_h}V_P\alpha_P^{n+1}
  =
  \sum_{P\in\Omega_h}V_P\alpha_P^n
  -
  \Delta t
  \sum_{f\in\partial\Omega_h}
  \bar{F}_{\alpha,f}.
\end{equation}

Internal face fluxes cancel pairwise.  Hence alpha changes only through the
boundary of the selected cell set.

The limiter is designed so that every explicit correction obeys the local
admissible bounds:

\begin{equation}
  \alpha_P^{lo}
  \le
  \alpha_P^{new}
  \le
  \alpha_P^{hi},
\end{equation}

subject to the usual finite-precision and linear-solver tolerances.  Since the
global hard bounds are included in the local interval, this also enforces

\begin{equation}
  0\le\alpha_P^{new}\le 1.
\end{equation}

The density conservation relation \eqref{eq:density-consistency} then connects
bounded alpha transport to momentum advection through the published
\(\bar{F}_{\rho,f}\).

\section{Relation to the interFoam Pattern}

The MOOSE formulation and OpenFOAM's \texttt{interFoam} share the same broad
VOF/PIMPLE structure:

\begin{itemize}[leftmargin=2.5em]
  \item one velocity and one pressure field for two immiscible incompressible
        phases;
  \item volume-fraction transport with artificial compression;
  \item bounded flux correction for alpha;
  \item mixture density and viscosity computed from alpha;
  \item pressure-velocity coupling through PIMPLE;
  \item reduced pressure \(p_{rgh}\) for hydrostatic effects.
\end{itemize}

The MOOSE implementation described here is specialized around an
alpha-consistent mass flux and a sharp-interface Rhie-Chow object.  The key
implementation distinction is that the same limited alpha flux that updates
\(\alpha\) also defines the density flux consumed by momentum.  This is a
conservative coupling contract between the free-surface submodule and the
Navier-Stokes finite-volume kernels.

\section{References}

\begin{thebibliography}{9}

\bibitem{openfoam-interfoam}
OpenCFD Ltd.,
\emph{OpenFOAM User Guide: interFoam}.
\url{https://www.openfoam.com/documentation/guides/latest/doc/guide-applications-solvers-multiphase-interFoam.html}

\bibitem{moose-dam-break-input}
MOOSE local input:
\texttt{dam\_break\_openfoam\_geometry.i}.

\bibitem{moose-algorithm-note}
MOOSE local documentation:
\texttt{modules/navier\_stokes/doc/content/modules/navier\_stokes/sharp\_interface\_dam\_break\_algorithm.md}.

\bibitem{moose-vof-corrector}
MOOSE local source:
\texttt{modules/navier\_stokes/src/userobjects/ConservativeSharpInterfaceVOFMULESCorrector.C}.

\bibitem{moose-rhiechow}
MOOSE local source:
\texttt{modules/navier\_stokes/src/userobjects/ConservativeSharpInterfaceRhieChowMassFluxBase.C}.

\bibitem{moose-executioner}
MOOSE local source:
\texttt{modules/navier\_stokes/src/executioners/ReducedPressurePIMPLESolve.C}.

\end{thebibliography}

\end{document}
